\Chapter{114}

\Verse{1}
{
    \zR{却说宝玉宝钗听说凤姐病的危急，赶忙起来，丫头秉烛伺候。}
}
{
    \eR{[English source not present in the checked-in Bencraft/Joly files.]}
}
{
}

\Verse{2}
{
    \zR{正要出院，只见 王夫人那边打发人来说：“琏二奶奶不好了，还没有咽气，二爷二奶奶且慢些过去 罢。}
}
{
    \eR{[English source not present in the checked-in Bencraft/Joly files.]}
}
{
}

\Verse{3}
{
    \zR{琏二奶奶的病有些古怪，从三更天起，到四更时候，没有住嘴，说了好些胡话， 要船要轿，只说赶到金陵归入什么册子去。}
}
{
    \eR{[English source not present in the checked-in Bencraft/Joly files.]}
}
{
}

\Verse{4}
{
    \zR{众人不懂，他只是哭哭喊喊。}
}
{
    \eR{[English source not present in the checked-in Bencraft/Joly files.]}
}
{
}

\Verse{5}
{
    \zR{琏二爷没 有法儿，只得去糊船轿，还没拿来，琏二奶奶喘着气等着呢。}
}
{
    \eR{[English source not present in the checked-in Bencraft/Joly files.]}
}
{
}

\Verse{6}
{
    \zR{太太叫我们过来说， 等琏二奶奶去了，再过去罢。”}
}
{
    \eR{[English source not present in the checked-in Bencraft/Joly files.]}
}
{
}

\Verse{7}
{
    \zR{宝玉道：“这也奇，他到金陵做什么去？”}
}
{
    \eR{[English source not present in the checked-in Bencraft/Joly files.]}
}
{
}

\Verse{8}
{
    \zR{袭人轻 轻的说道：“你不是那年做梦，我还记得说有多少册子？}
}
{
    \eR{[English source not present in the checked-in Bencraft/Joly files.]}
}
{
}

\Verse{9}
{
    \zR{莫不琏二奶奶是到那里去 罢？”}
}
{
    \eR{[English source not present in the checked-in Bencraft/Joly files.]}
}
{
}

\Verse{10}
{
    \zR{宝玉听了点头道：“是呀，可惜我都不记得那上头的话了。}
}
{
    \eR{[English source not present in the checked-in Bencraft/Joly files.]}
}
{
}

\Verse{11}
{
    \zR{这么说起来，人 都有个定数的了。}
}
{
    \eR{[English source not present in the checked-in Bencraft/Joly files.]}
}
{
}

\Verse{12}
{
    \zR{但不知林妹妹又到那里去了？}
}
{
    \eR{[English source not present in the checked-in Bencraft/Joly files.]}
}
{
}

\Verse{13}
{
    \zR{我如今被你一说，我有些懂的了。}
}
{
    \eR{[English source not present in the checked-in Bencraft/Joly files.]}
}
{
}

\Verse{14}
{
    \zR{若再做这个梦时，我必细细的瞧一瞧，便有未卜先知的分儿了。”}
}
{
    \eR{[English source not present in the checked-in Bencraft/Joly files.]}
}
{
}

\Verse{15}
{
    \zR{袭人道：“你这 样的人，可是不可合你说话，我偶然提了一句，你就认起真来了吗？}
}
{
    \eR{[English source not present in the checked-in Bencraft/Joly files.]}
}
{
}

\Verse{16}
{
    \zR{就算你能先知 了，又有什么法儿？”}
}
{
    \eR{[English source not present in the checked-in Bencraft/Joly files.]}
}
{
}

\Verse{17}
{
    \zR{宝玉道：“只怕不能先知；若是能了，我也犯不着为你们瞎 操心了。”}
}
{
    \eR{[English source not present in the checked-in Bencraft/Joly files.]}
}
{
}

\Verse{18}
{
    \zR{两人正说着，宝钗走来，问道：“你们说什么？”}
}
{
    \eR{[English source not present in the checked-in Bencraft/Joly files.]}
}
{
}

\Verse{19}
{
    \zR{宝玉恐他盘诘，只说： “我们谈论凤姐姐。”}
}
{
    \eR{[English source not present in the checked-in Bencraft/Joly files.]}
}
{
}

\Verse{20}
{
    \zR{宝钗道：“人要死了，你们还只管议论他。}
}
{
    \eR{[English source not present in the checked-in Bencraft/Joly files.]}
}
{
}

\Verse{21}
{
    \zR{旧年你还说我咒 人，那个签不是应了么？”}
}
{
    \eR{[English source not present in the checked-in Bencraft/Joly files.]}
}
{
}

\Verse{22}
{
    \zR{宝玉又想了一想，拍手道：“是的是的，这么说起来， 你倒能先知了。}
}
{
    \eR{[English source not present in the checked-in Bencraft/Joly files.]}
}
{
}

\Verse{23}
{
    \zR{我索性问问你，你知道我将来怎么样？”}
}
{
    \eR{[English source not present in the checked-in Bencraft/Joly files.]}
}
{
}

\Verse{24}
{
    \zR{宝钗笑道：“这是又胡闹 起来了。}
}
{
    \eR{[English source not present in the checked-in Bencraft/Joly files.]}
}
{
}

\Verse{25}
{
    \zR{我是就他求的签上的话混解的，你就认了真了。}
}
{
    \eR{[English source not present in the checked-in Bencraft/Joly files.]}
}
{
}

\Verse{26}
{
    \zR{你和我们二嫂子成了一样 的了。}
}
{
    \eR{[English source not present in the checked-in Bencraft/Joly files.]}
}
{
}

\Verse{27}
{
    \zR{你失了玉，他去求妙玉扶乩，批出来众人不解。}
}
{
    \eR{[English source not present in the checked-in Bencraft/Joly files.]}
}
{
}

\Verse{28}
{
    \zR{他背地里合我说，妙玉怎么 前知，怎么参禅悟道，如今他遭此大难，如何自己都不知道？}
}
{
    \eR{[English source not present in the checked-in Bencraft/Joly files.]}
}
{
}

\Verse{29}
{
    \zR{这可是算得前知吗？}
}
{
    \eR{[English source not present in the checked-in Bencraft/Joly files.]}
}
{
}

\Verse{30}
{
    \zR{就 是我偶然说着了二奶奶的事情，其实知道他是怎么样了？}
}
{
    \eR{[English source not present in the checked-in Bencraft/Joly files.]}
}
{
}

\Verse{31}
{
    \zR{只怕我连我自己也不知道 呢。}
}
{
    \eR{[English source not present in the checked-in Bencraft/Joly files.]}
}
{
}

\Verse{32}
{
    \zR{这些事情，原都是虚诞的，可是信得的么？”}
}
{
    \eR{[English source not present in the checked-in Bencraft/Joly files.]}
}
{
}

\Verse{33}
{
    \zR{宝玉道：“别提他了。}
}
{
    \eR{[English source not present in the checked-in Bencraft/Joly files.]}
}
{
}

\Verse{34}
{
    \zR{你只说邢妹妹罢，自从我们这里连连的有事，把他这件 事竟忘记了。}
}
{
    \eR{[English source not present in the checked-in Bencraft/Joly files.]}
}
{
}

\Verse{35}
{
    \zR{你们家这么一件大事，怎么就草草的完了？}
}
{
    \eR{[English source not present in the checked-in Bencraft/Joly files.]}
}
{
}

\Verse{36}
{
    \zR{也没请亲唤友的。”}
}
{
    \eR{[English source not present in the checked-in Bencraft/Joly files.]}
}
{
}

\Verse{37}
{
    \zR{宝钗 道：“你这话又是迂了。}
}
{
    \eR{[English source not present in the checked-in Bencraft/Joly files.]}
}
{
}

\Verse{38}
{
    \zR{我们家的亲戚，只有咱们这里和王家最近。}
}
{
    \eR{[English source not present in the checked-in Bencraft/Joly files.]}
}
{
}

\Verse{39}
{
    \zR{王家没了什么 正经人了，咱们家遭了老太太的大事，所以也没请，就是琏二哥张罗了张罗。}
}
{
    \eR{[English source not present in the checked-in Bencraft/Joly files.]}
}
{
}

\Verse{40}
{
    \zR{别的 亲戚虽也有一两门子，你没过去，如何知道？}
}
{
    \eR{[English source not present in the checked-in Bencraft/Joly files.]}
}
{
}

\Verse{41}
{
    \zR{算起来，我们这二嫂子的命和我差不 多。}
}
{
    \eR{[English source not present in the checked-in Bencraft/Joly files.]}
}
{
}

\Verse{42}
{
    \zR{好好的许了我二哥哥，我妈妈原想要体体面面的给二哥哥娶这房亲事的。}
}
{
    \eR{[English source not present in the checked-in Bencraft/Joly files.]}
}
{
}

\Verse{43}
{
    \zR{一则 为我哥哥在监里，二哥哥也不肯大办；二则为咱们家的事；三则为我二嫂子在大太 太那边忒苦，又加着抄了家，大太太是一味的苛刻，他也实在难受。}
}
{
    \eR{[English source not present in the checked-in Bencraft/Joly files.]}
}
{
}

\Verse{44}
{
    \zR{所以我和妈妈 说了，便将将就就的娶了过去。}
}
{
    \eR{[English source not present in the checked-in Bencraft/Joly files.]}
}
{
}

\Verse{45}
{
    \zR{我看二嫂子如今倒是安心乐意的孝敬我妈妈，比亲 媳妇还强十倍呢。}
}
{
    \eR{[English source not present in the checked-in Bencraft/Joly files.]}
}
{
}

\Verse{46}
{
    \zR{待二哥哥也是极尽妇道的，和香菱又甚好。}
}
{
    \eR{[English source not present in the checked-in Bencraft/Joly files.]}
}
{
}

\Verse{47}
{
    \zR{二哥哥不在家，他两 个和和气气的过日子，虽说是穷些，我妈妈近来倒安逸好些。}
}
{
    \eR{[English source not present in the checked-in Bencraft/Joly files.]}
}
{
}

\Verse{48}
{
    \zR{就是想起我哥哥来， 不免伤心。}
}
{
    \eR{[English source not present in the checked-in Bencraft/Joly files.]}
}
{
}

\Verse{49}
{
    \zR{况且常打发人家里来要使用，多亏二哥哥在外头账头儿上讨来应付他。}
}
{
    \eR{[English source not present in the checked-in Bencraft/Joly files.]}
}
{
}

\Verse{50}
{
    \zR{我听见说：城里的几处房子已经也典了，还剩了一所，如今打算着搬了去住。”}
}
{
    \eR{[English source not present in the checked-in Bencraft/Joly files.]}
}
{
}

\Verse{51}
{
    \zR{宝 玉道：“为什么要搬？}
}
{
    \eR{[English source not present in the checked-in Bencraft/Joly files.]}
}
{
}

\Verse{52}
{
    \zR{住在这里，你来去也便宜些；若搬远了，你去就要一天了。”}
}
{
    \eR{[English source not present in the checked-in Bencraft/Joly files.]}
}
{
}

\Verse{53}
{
    \zR{宝钗道：“虽说是亲戚，到底各自的稳便些。}
}
{
    \eR{[English source not present in the checked-in Bencraft/Joly files.]}
}
{
}

\Verse{54}
{
    \zR{那里有个一辈子住在亲戚家的呢？”}
}
{
    \eR{[English source not present in the checked-in Bencraft/Joly files.]}
}
{
}

\Verse{55}
{
    \zR{宝玉还要讲出不搬去的理，王夫人打发人来说：“琏二奶奶咽了气了!所有的 人都过去了，请二爷二奶奶就过去。”}
}
{
    \eR{[English source not present in the checked-in Bencraft/Joly files.]}
}
{
}

\Verse{56}
{
    \zR{宝玉听了，也掌不住跺脚要哭。}
}
{
    \eR{[English source not present in the checked-in Bencraft/Joly files.]}
}
{
}

\Verse{57}
{
    \zR{宝钗虽也悲 戚，恐宝玉伤心，便说：“有在这里哭的，不如到那边哭去。”}
}
{
    \eR{[English source not present in the checked-in Bencraft/Joly files.]}
}
{
}

\Verse{58}
{
    \zR{于是两人一直到凤 姐那里，只见好些人围着哭呢。}
}
{
    \eR{[English source not present in the checked-in Bencraft/Joly files.]}
}
{
}

\Verse{59}
{
    \zR{宝钗走到跟前，见凤姐已经停床，便大放悲声。}
}
{
    \eR{[English source not present in the checked-in Bencraft/Joly files.]}
}
{
}

\Verse{60}
{
    \zR{宝 玉也拉着贾琏的手，大哭起来，贾琏也重新哭泣。}
}
{
    \eR{[English source not present in the checked-in Bencraft/Joly files.]}
}
{
}

\Verse{61}
{
    \zR{平儿等因见无人劝解，只得含悲 上来劝止了。}
}
{
    \eR{[English source not present in the checked-in Bencraft/Joly files.]}
}
{
}

\Verse{62}
{
    \zR{众人都悲哀不止。}
}
{
    \eR{[English source not present in the checked-in Bencraft/Joly files.]}
}
{
}

\Verse{63}
{
    \zR{贾琏此时手足无措，叫人传了赖大来，叫他办理丧 事。}
}
{
    \eR{[English source not present in the checked-in Bencraft/Joly files.]}
}
{
}

\Verse{64}
{
    \zR{自己回明了贾政，然后去行事。}
}
{
    \eR{[English source not present in the checked-in Bencraft/Joly files.]}
}
{
}

\Verse{65}
{
    \zR{但是手头不济，诸事拮据。}
}
{
    \eR{[English source not present in the checked-in Bencraft/Joly files.]}
}
{
}

\Verse{66}
{
    \zR{又想起凤姐素日的 好处来，更加悲哭不已。}
}
{
    \eR{[English source not present in the checked-in Bencraft/Joly files.]}
}
{
}

\Verse{67}
{
    \zR{又见巧姐哭的死去活来，越发伤心。}
}
{
    \eR{[English source not present in the checked-in Bencraft/Joly files.]}
}
{
}

\Verse{68}
{
    \zR{哭到天明，即刻打发 人去请他大舅子王仁过来。}
}
{
    \eR{[English source not present in the checked-in Bencraft/Joly files.]}
}
{
}

\Verse{69}
{
    \zR{那王仁自从王子腾死后，王子胜又是无能的人，任他胡为，已闹的六亲不和。}
}
{
    \eR{[English source not present in the checked-in Bencraft/Joly files.]}
}
{
}

\Verse{70}
{
    \zR{今知妹子死了，只得赶着过来哭了一场。}
}
{
    \eR{[English source not present in the checked-in Bencraft/Joly files.]}
}
{
}

\Verse{71}
{
    \zR{见这里诸事将就，心下便不舒服，说：“我 妹妹在你家辛辛苦苦当了好几年家，也没有什么错处，你们家该认真的发送发送才 是，怎么这时候诸事还没有齐备？”}
}
{
    \eR{[English source not present in the checked-in Bencraft/Joly files.]}
}
{
}

\Verse{72}
{
    \zR{贾琏本与王仁不睦，见他说些混帐话，知他不 懂的什么，也不大理他。}
}
{
    \eR{[English source not present in the checked-in Bencraft/Joly files.]}
}
{
}

\Verse{73}
{
    \zR{王仁便叫了他外甥女儿巧姐过来，说：“你娘在时，本来 办事不周到：只知道一味的奉承老太太，把我们的人都不大看在眼里。}
}
{
    \eR{[English source not present in the checked-in Bencraft/Joly files.]}
}
{
}

\Verse{74}
{
    \zR{外甥女儿! 你也大了，看见我从来沾染过你们没有？}
}
{
    \eR{[English source not present in the checked-in Bencraft/Joly files.]}
}
{
}

\Verse{75}
{
    \zR{如今你娘死了，诸事要听着舅舅的话。}
}
{
    \eR{[English source not present in the checked-in Bencraft/Joly files.]}
}
{
}

\Verse{76}
{
    \zR{你 母亲娘家的亲戚就是我和你二舅舅了。}
}
{
    \eR{[English source not present in the checked-in Bencraft/Joly files.]}
}
{
}

\Verse{77}
{
    \zR{你父亲的为人，我也早知道了，只有敬重别 人的。}
}
{
    \eR{[English source not present in the checked-in Bencraft/Joly files.]}
}
{
}

\Verse{78}
{
    \zR{那年什么尤姨娘死了，我虽不在京，听见说花了好些银子。}
}
{
    \eR{[English source not present in the checked-in Bencraft/Joly files.]}
}
{
}

\Verse{79}
{
    \zR{如今你娘死了， 你父亲倒是这样的将就办去，你也不知道劝劝你父亲吗？”}
}
{
    \eR{[English source not present in the checked-in Bencraft/Joly files.]}
}
{
}

\Verse{80}
{
    \zR{巧姐道：“我父亲巴不 得要好看，只是如今比不得从前了。}
}
{
    \eR{[English source not present in the checked-in Bencraft/Joly files.]}
}
{
}

\Verse{81}
{
    \zR{现在手里没钱，所以诸事省些是有的。”}
}
{
    \eR{[English source not present in the checked-in Bencraft/Joly files.]}
}
{
}

\Verse{82}
{
    \zR{王仁 道：“你的东西还少么？”}
}
{
    \eR{[English source not present in the checked-in Bencraft/Joly files.]}
}
{
}

\Verse{83}
{
    \zR{巧姐儿道：“旧年抄去，何尝还有呢？”}
}
{
    \eR{[English source not present in the checked-in Bencraft/Joly files.]}
}
{
}

\Verse{84}
{
    \zR{王仁道：“你 也这样说？}
}
{
    \eR{[English source not present in the checked-in Bencraft/Joly files.]}
}
{
}

\Verse{85}
{
    \zR{我听见老太太又给了好些东西，你该拿出来。”}
}
{
    \eR{[English source not present in the checked-in Bencraft/Joly files.]}
}
{
}

\Verse{86}
{
    \zR{巧姐又不好说父亲用去， 只推不知道。}
}
{
    \eR{[English source not present in the checked-in Bencraft/Joly files.]}
}
{
}

\Verse{87}
{
    \zR{王仁便道：“哦，我知道了，不过是你要留着做嫁妆罢咧。”}
}
{
    \eR{[English source not present in the checked-in Bencraft/Joly files.]}
}
{
}

\Verse{88}
{
    \zR{巧姐听 了，不敢回言，只气得哽噎难鸣的哭起来了。}
}
{
    \eR{[English source not present in the checked-in Bencraft/Joly files.]}
}
{
}

\Verse{89}
{
    \zR{平儿生气说道：“舅老爷，有话等我 们二爷进来再说。}
}
{
    \eR{[English source not present in the checked-in Bencraft/Joly files.]}
}
{
}

\Verse{90}
{
    \zR{姑娘这么点年纪，他懂的什么？”}
}
{
    \eR{[English source not present in the checked-in Bencraft/Joly files.]}
}
{
}

\Verse{91}
{
    \zR{王仁道：“你们是巴不得二奶 奶死了，你们就好为王了。}
}
{
    \eR{[English source not present in the checked-in Bencraft/Joly files.]}
}
{
}

\Verse{92}
{
    \zR{我并不要什么，好看些，也是你们的脸面！”}
}
{
    \eR{[English source not present in the checked-in Bencraft/Joly files.]}
}
{
}

\Verse{93}
{
    \zR{说着赌气 坐着。}
}
{
    \eR{[English source not present in the checked-in Bencraft/Joly files.]}
}
{
}

\Verse{94}
{
    \zR{巧姐满心的不舒服，心想：“我父亲并不是没情。}
}
{
    \eR{[English source not present in the checked-in Bencraft/Joly files.]}
}
{
}

\Verse{95}
{
    \zR{我妈妈在时，舅舅不知拿 了多少东西去，如今说得这样干净！”}
}
{
    \eR{[English source not present in the checked-in Bencraft/Joly files.]}
}
{
}

\Verse{96}
{
    \zR{于是便不大瞧得起他舅舅了。}
}
{
    \eR{[English source not present in the checked-in Bencraft/Joly files.]}
}
{
}

\Verse{97}
{
    \zR{岂知王仁心里 想来，他妹妹不知积攒了多少。}
}
{
    \eR{[English source not present in the checked-in Bencraft/Joly files.]}
}
{
}

\Verse{98}
{
    \zR{“虽说抄了家，那屋里的银子还怕少吗？}
}
{
    \eR{[English source not present in the checked-in Bencraft/Joly files.]}
}
{
}

\Verse{99}
{
    \zR{必是怕我 来缠他们，所以也帮着这么说。}
}
{
    \eR{[English source not present in the checked-in Bencraft/Joly files.]}
}
{
}

\Verse{100}
{
    \zR{这小东西儿也是不中用的！”}
}
{
    \eR{[English source not present in the checked-in Bencraft/Joly files.]}
}
{
}

\Verse{101}
{
    \zR{从此王仁也嫌了巧姐 儿了。}
}
{
    \eR{[English source not present in the checked-in Bencraft/Joly files.]}
}
{
}

\Verse{102}
{
    \zR{贾琏并不知道，只忙着弄银钱使用。}
}
{
    \eR{[English source not present in the checked-in Bencraft/Joly files.]}
}
{
}

\Verse{103}
{
    \zR{外头的大事叫赖大办了，里头也要用好些 钱，一时实在不能张罗。}
}
{
    \eR{[English source not present in the checked-in Bencraft/Joly files.]}
}
{
}

\Verse{104}
{
    \zR{平儿知他着急，便叫贾琏道：“二爷也别过于伤了自己的 身子。”}
}
{
    \eR{[English source not present in the checked-in Bencraft/Joly files.]}
}
{
}

\Verse{105}
{
    \zR{贾琏道：“什么身子!现在日用的钱都没有，这件事怎么办？}
}
{
    \eR{[English source not present in the checked-in Bencraft/Joly files.]}
}
{
}

\Verse{106}
{
    \zR{偏有个糊涂行 子又在这里蛮缠，你想有什么法儿？”}
}
{
    \eR{[English source not present in the checked-in Bencraft/Joly files.]}
}
{
}

\Verse{107}
{
    \zR{平儿道：“二爷也不用着急。}
}
{
    \eR{[English source not present in the checked-in Bencraft/Joly files.]}
}
{
}

\Verse{108}
{
    \zR{若说没钱使唤， 我还有些东西，旧年幸亏没有抄在里头去，二爷要，就拿去当着使唤罢。”}
}
{
    \eR{[English source not present in the checked-in Bencraft/Joly files.]}
}
{
}

\Verse{109}
{
    \zR{贾琏听 了，心想：“难得这样。”}
}
{
    \eR{[English source not present in the checked-in Bencraft/Joly files.]}
}
{
}

\Verse{110}
{
    \zR{便笑道：“这样更好，省得我各处张罗。}
}
{
    \eR{[English source not present in the checked-in Bencraft/Joly files.]}
}
{
}

\Verse{111}
{
    \zR{等我银子弄到 手了还你。”}
}
{
    \eR{[English source not present in the checked-in Bencraft/Joly files.]}
}
{
}

\Verse{112}
{
    \zR{平儿道：“我的也是奶奶给的，什么还不还。}
}
{
    \eR{[English source not present in the checked-in Bencraft/Joly files.]}
}
{
}

\Verse{113}
{
    \zR{只要这件事办的好看些 就是了。”}
}
{
    \eR{[English source not present in the checked-in Bencraft/Joly files.]}
}
{
}

\Verse{114}
{
    \zR{贾琏心里倒着实感激他，便将平儿的东西拿了去，当钱使用。}
}
{
    \eR{[English source not present in the checked-in Bencraft/Joly files.]}
}
{
}

\Verse{115}
{
    \zR{诸凡事情， 便与平儿商量。}
}
{
    \eR{[English source not present in the checked-in Bencraft/Joly files.]}
}
{
}

\Verse{116}
{
    \zR{秋桐看着，心里就有些不甘，每每口角里头便说：“平儿没有了奶 奶，他要上去了。}
}
{
    \eR{[English source not present in the checked-in Bencraft/Joly files.]}
}
{
}

\Verse{117}
{
    \zR{我是老爷的人，他怎么就越过我去了呢？”}
}
{
    \eR{[English source not present in the checked-in Bencraft/Joly files.]}
}
{
}

\Verse{118}
{
    \zR{平儿也看出来了，只 不理他。}
}
{
    \eR{[English source not present in the checked-in Bencraft/Joly files.]}
}
{
}

\Verse{119}
{
    \zR{倒是贾琏一时明白，越发把秋桐嫌了，碰着有些烦恼，便拿着秋桐出气。}
}
{
    \eR{[English source not present in the checked-in Bencraft/Joly files.]}
}
{
}

\Verse{120}
{
    \zR{邢夫人知道，反说贾琏不好。}
}
{
    \eR{[English source not present in the checked-in Bencraft/Joly files.]}
}
{
}

\Verse{121}
{
    \zR{贾琏忍气不提。}
}
{
    \eR{[English source not present in the checked-in Bencraft/Joly files.]}
}
{
}

\Verse{122}
{
    \zR{再说凤姐停了十馀天，送了殡。}
}
{
    \eR{[English source not present in the checked-in Bencraft/Joly files.]}
}
{
}

\Verse{123}
{
    \zR{贾政守着老太太的孝，总在外书房。}
}
{
    \eR{[English source not present in the checked-in Bencraft/Joly files.]}
}
{
}

\Verse{124}
{
    \zR{那时清客 相公，渐渐的都辞去了，只有个程日兴还在那里，时常陪着说说话儿。}
}
{
    \eR{[English source not present in the checked-in Bencraft/Joly files.]}
}
{
}

\Verse{125}
{
    \zR{提起：“家 运不好，一连人口死了好些，大老爷合珍大爷又在外头。}
}
{
    \eR{[English source not present in the checked-in Bencraft/Joly files.]}
}
{
}

\Verse{126}
{
    \zR{家计一天难似一天，外头 东庄地亩也不知道怎么样，总不得了！”}
}
{
    \eR{[English source not present in the checked-in Bencraft/Joly files.]}
}
{
}

\Verse{127}
{
    \zR{那程日兴道：“我在这里好些年，也知道 府上的人，那一个不是肥己的？}
}
{
    \eR{[English source not present in the checked-in Bencraft/Joly files.]}
}
{
}

\Verse{128}
{
    \zR{一年一年都往他家里拿，那自然府上是一年不够一 年了。}
}
{
    \eR{[English source not present in the checked-in Bencraft/Joly files.]}
}
{
}

\Verse{129}
{
    \zR{又添了大老爷珍大爷那边两处的费用，外头又有些债务。}
}
{
    \eR{[English source not present in the checked-in Bencraft/Joly files.]}
}
{
}

\Verse{130}
{
    \zR{前儿又破了好些财， 要想衙门里缉贼追赃，那是难事。}
}
{
    \eR{[English source not present in the checked-in Bencraft/Joly files.]}
}
{
}

\Verse{131}
{
    \zR{老世翁若要安顿家事，除非传那些管事的来，派 一个心腹人各处去清查清查：该去的去，该留的留；有了亏空，着在经手的身上赔 补，这就有了数儿了。}
}
{
    \eR{[English source not present in the checked-in Bencraft/Joly files.]}
}
{
}

\Verse{132}
{
    \zR{那一座大园子，人家是不敢买的，这里头的出息也不少，又 不派人管了。}
}
{
    \eR{[English source not present in the checked-in Bencraft/Joly files.]}
}
{
}

\Verse{133}
{
    \zR{几年老世翁不在家，这些人就弄神弄鬼儿的，闹的一个人不敢到园里， 这都是家人的弊。}
}
{
    \eR{[English source not present in the checked-in Bencraft/Joly files.]}
}
{
}

\Verse{134}
{
    \zR{此时把下人查一查，好的使着，不好的便撵了，这才是道理。”}
}
{
    \eR{[English source not present in the checked-in Bencraft/Joly files.]}
}
{
}

\Verse{135}
{
    \zR{贾政点头道：“先生你有所不知!不必说下人，就是自己的侄儿，也靠不住!若要我 查起来，那能一一亲见亲知？}
}
{
    \eR{[English source not present in the checked-in Bencraft/Joly files.]}
}
{
}

\Verse{136}
{
    \zR{况我又在服中，不能照管这些个。}
}
{
    \eR{[English source not present in the checked-in Bencraft/Joly files.]}
}
{
}

\Verse{137}
{
    \zR{我素来又兼不大理 家，有的没的，我还摸不着呢。”}
}
{
    \eR{[English source not present in the checked-in Bencraft/Joly files.]}
}
{
}

\Verse{138}
{
    \zR{程日兴道：“老世翁最是仁德的人。}
}
{
    \eR{[English source not present in the checked-in Bencraft/Joly files.]}
}
{
}

\Verse{139}
{
    \zR{若在别人家 这样的家计，就穷起来，十年五载还不怕，便向这些管家的要，也就够了。}
}
{
    \eR{[English source not present in the checked-in Bencraft/Joly files.]}
}
{
}

\Verse{140}
{
    \zR{我听见 世翁的家人还有做知县的呢。”}
}
{
    \eR{[English source not present in the checked-in Bencraft/Joly files.]}
}
{
}

\Verse{141}
{
    \zR{贾政道：“一个人若要使起家人们的钱来，便了不 得了，只好自己俭省些。}
}
{
    \eR{[English source not present in the checked-in Bencraft/Joly files.]}
}
{
}

\Verse{142}
{
    \zR{但是册子上的产业，若是实有还好，生怕有名无实了。”}
}
{
    \eR{[English source not present in the checked-in Bencraft/Joly files.]}
}
{
}

\Verse{143}
{
    \zR{程日兴道：“老世翁所见极是。}
}
{
    \eR{[English source not present in the checked-in Bencraft/Joly files.]}
}
{
}

\Verse{144}
{
    \zR{晚生为什么说要查查呢！”}
}
{
    \eR{[English source not present in the checked-in Bencraft/Joly files.]}
}
{
}

\Verse{145}
{
    \zR{贾政道：“先生必有所 闻？”}
}
{
    \eR{[English source not present in the checked-in Bencraft/Joly files.]}
}
{
}

\Verse{146}
{
    \zR{程日兴道：“我虽知道些那些管事的神通，晚生也不敢言语的。”}
}
{
    \eR{[English source not present in the checked-in Bencraft/Joly files.]}
}
{
}

\Verse{147}
{
    \zR{贾政听了， 便知话里有因，便叹道：“我家祖父以来，都是仁厚的，从没有刻薄过下人。}
}
{
    \eR{[English source not present in the checked-in Bencraft/Joly files.]}
}
{
}

\Verse{148}
{
    \zR{我看 如今这些人一日不似一日了。}
}
{
    \eR{[English source not present in the checked-in Bencraft/Joly files.]}
}
{
}

\Verse{149}
{
    \zR{在我手里行出主子样儿来，又叫人笑话。”}
}
{
    \eR{[English source not present in the checked-in Bencraft/Joly files.]}
}
{
}

\Verse{150}
{
    \zR{两人正说着，门上的进来回道：“江南甄老爷来了。”}
}
{
    \eR{[English source not present in the checked-in Bencraft/Joly files.]}
}
{
}

\Verse{151}
{
    \zR{贾政便问道：“甄老爷 进京为什么？”}
}
{
    \eR{[English source not present in the checked-in Bencraft/Joly files.]}
}
{
}

\Verse{152}
{
    \zR{那人道：“奴才也打听过了，说是蒙圣恩起复了。”}
}
{
    \eR{[English source not present in the checked-in Bencraft/Joly files.]}
}
{
}

\Verse{153}
{
    \zR{贾政道：“不 用说了，快请罢。”}
}
{
    \eR{[English source not present in the checked-in Bencraft/Joly files.]}
}
{
}

\Verse{154}
{
    \zR{那人出去，请了进来。}
}
{
    \eR{[English source not present in the checked-in Bencraft/Joly files.]}
}
{
}

\Verse{155}
{
    \zR{那甄老爷即是甄宝玉之父，名叫甄应嘉， 表字友忠，也是金陵人氏，功勋之后。}
}
{
    \eR{[English source not present in the checked-in Bencraft/Joly files.]}
}
{
}

\Verse{156}
{
    \zR{原与贾府有亲，素来走动的。}
}
{
    \eR{[English source not present in the checked-in Bencraft/Joly files.]}
}
{
}

\Verse{157}
{
    \zR{因前年挂误革 了职，动了家产，今遇主上眷念功臣，赐还世职，行取来京陛见。}
}
{
    \eR{[English source not present in the checked-in Bencraft/Joly files.]}
}
{
}

\Verse{158}
{
    \zR{知道贾母新丧， 特备祭礼，择日到寄灵的地方拜奠，所以先来拜望。}
}
{
    \eR{[English source not present in the checked-in Bencraft/Joly files.]}
}
{
}

\Verse{159}
{
    \zR{贾政有服，不能远接，在外书房门口等着。}
}
{
    \eR{[English source not present in the checked-in Bencraft/Joly files.]}
}
{
}

\Verse{160}
{
    \zR{那位甄老爷一见，便悲喜交集。}
}
{
    \eR{[English source not present in the checked-in Bencraft/Joly files.]}
}
{
}

\Verse{161}
{
    \zR{因 在制中，不便行礼，遂拉着手叙了些阔别思念的话。}
}
{
    \eR{[English source not present in the checked-in Bencraft/Joly files.]}
}
{
}

\Verse{162}
{
    \zR{然后分宾主坐下，献了茶，彼 此又将别后事情的话说了。}
}
{
    \eR{[English source not present in the checked-in Bencraft/Joly files.]}
}
{
}

\Verse{163}
{
    \zR{贾政问道：“老亲翁几时陛见的？”}
}
{
    \eR{[English source not present in the checked-in Bencraft/Joly files.]}
}
{
}

\Verse{164}
{
    \zR{甄应嘉道：“前日。”}
}
{
    \eR{[English source not present in the checked-in Bencraft/Joly files.]}
}
{
}

\Verse{165}
{
    \zR{贾政道：“主上隆恩，必有温谕。”}
}
{
    \eR{[English source not present in the checked-in Bencraft/Joly files.]}
}
{
}

\Verse{166}
{
    \zR{甄应嘉道：“主上的恩典，真是比天还高，下 了好些旨意。”}
}
{
    \eR{[English source not present in the checked-in Bencraft/Joly files.]}
}
{
}

\Verse{167}
{
    \zR{贾政道：“什么好旨意？”}
}
{
    \eR{[English source not present in the checked-in Bencraft/Joly files.]}
}
{
}

\Verse{168}
{
    \zR{甄应嘉道：“近来越寇猖獗，海疆一带， 小民不安，派了安国公征剿贼寇。}
}
{
    \eR{[English source not present in the checked-in Bencraft/Joly files.]}
}
{
}

\Verse{169}
{
    \zR{主上因我熟悉土疆，命我前往安抚，但是即日就 要起身。}
}
{
    \eR{[English source not present in the checked-in Bencraft/Joly files.]}
}
{
}

\Verse{170}
{
    \zR{昨日知老太太仙逝，谨备瓣香至灵前拜奠，稍尽微忱。”}
}
{
    \eR{[English source not present in the checked-in Bencraft/Joly files.]}
}
{
}

\Verse{171}
{
    \zR{贾政即忙叩首拜 谢，便说：“老亲翁即此一行，必是上慰圣心，下安黎庶。}
}
{
    \eR{[English source not present in the checked-in Bencraft/Joly files.]}
}
{
}

\Verse{172}
{
    \zR{诚哉莫大之功，正在此 行。}
}
{
    \eR{[English source not present in the checked-in Bencraft/Joly files.]}
}
{
}

\Verse{173}
{
    \zR{但弟不克亲睹奇才，只好遥聆捷报。}
}
{
    \eR{[English source not present in the checked-in Bencraft/Joly files.]}
}
{
}

\Verse{174}
{
    \zR{现在镇海统制是弟舍亲，会时务望青照。”}
}
{
    \eR{[English source not present in the checked-in Bencraft/Joly files.]}
}
{
}

\Verse{175}
{
    \zR{甄应嘉道：“老亲翁与统制是什么亲戚？”}
}
{
    \eR{[English source not present in the checked-in Bencraft/Joly files.]}
}
{
}

\Verse{176}
{
    \zR{贾政道：“弟那年在江西粮道任时，将 小女许配与统制少君，结？}
}
{
    \eR{[English source not present in the checked-in Bencraft/Joly files.]}
}
{
}

\Verse{177}
{
    \zR{已经三载。}
}
{
    \eR{[English source not present in the checked-in Bencraft/Joly files.]}
}
{
}

\Verse{178}
{
    \zR{因海口案内未清，继以海寇聚奸，所以音信 不通。}
}
{
    \eR{[English source not present in the checked-in Bencraft/Joly files.]}
}
{
}

\Verse{179}
{
    \zR{弟深念小女，俟老亲翁安抚事竣后，拜恳便中一视。}
}
{
    \eR{[English source not present in the checked-in Bencraft/Joly files.]}
}
{
}

\Verse{180}
{
    \zR{弟即修字数行，烦尊纪 带去，便感激不尽了。”}
}
{
    \eR{[English source not present in the checked-in Bencraft/Joly files.]}
}
{
}

\Verse{181}
{
    \zR{甄应嘉道：“儿女之情，人所不免。}
}
{
    \eR{[English source not present in the checked-in Bencraft/Joly files.]}
}
{
}

\Verse{182}
{
    \zR{我正在有奉托老亲翁 的事。}
}
{
    \eR{[English source not present in the checked-in Bencraft/Joly files.]}
}
{
}

\Verse{183}
{
    \zR{昨蒙圣恩召取来京，因小儿年幼，家下乏人，将贱眷全带来京。}
}
{
    \eR{[English source not present in the checked-in Bencraft/Joly files.]}
}
{
}

\Verse{184}
{
    \zR{我因钦限迅 速，昼夜先行，贱眷在后缓行，到京尚需时日。}
}
{
    \eR{[English source not present in the checked-in Bencraft/Joly files.]}
}
{
}

\Verse{185}
{
    \zR{弟奉旨出京，不敢久留。}
}
{
    \eR{[English source not present in the checked-in Bencraft/Joly files.]}
}
{
}

\Verse{186}
{
    \zR{将来贱眷 到京，少不得要到尊府，定叫小犬叩见。}
}
{
    \eR{[English source not present in the checked-in Bencraft/Joly files.]}
}
{
}

\Verse{187}
{
    \zR{如可进教，遇有姻事可图之处，望乞留意 为感。”}
}
{
    \eR{[English source not present in the checked-in Bencraft/Joly files.]}
}
{
}

\Verse{188}
{
    \zR{贾政一一答应。}
}
{
    \eR{[English source not present in the checked-in Bencraft/Joly files.]}
}
{
}

\Verse{189}
{
    \zR{那甄应嘉又说了几句话，就要起身，说：“明日在城外再 见。”}
}
{
    \eR{[English source not present in the checked-in Bencraft/Joly files.]}
}
{
}

\Verse{190}
{
    \zR{贾政见他事忙，谅难再坐，只得送出书房。}
}
{
    \eR{[English source not present in the checked-in Bencraft/Joly files.]}
}
{
}

\Verse{191}
{
    \zR{贾琏宝玉早已伺候在那里代送，因贾政未叫，不敢擅入。}
}
{
    \eR{[English source not present in the checked-in Bencraft/Joly files.]}
}
{
}

\Verse{192}
{
    \zR{甄应嘉出来，两人上 去请安。}
}
{
    \eR{[English source not present in the checked-in Bencraft/Joly files.]}
}
{
}

\Verse{193}
{
    \zR{应嘉一见宝玉，呆了一呆，心想：“这个怎么甚像我家宝玉!只是浑身缟 素。”}
}
{
    \eR{[English source not present in the checked-in Bencraft/Joly files.]}
}
{
}

\Verse{194}
{
    \zR{问道：“至亲久阔，爷们都不认得了。”}
}
{
    \eR{[English source not present in the checked-in Bencraft/Joly files.]}
}
{
}

\Verse{195}
{
    \zR{贾政忙指贾琏道：“这是家兄名赦 之子琏二侄儿。”}
}
{
    \eR{[English source not present in the checked-in Bencraft/Joly files.]}
}
{
}

\Verse{196}
{
    \zR{又指着宝玉道：“这是第二小犬，名叫宝玉。”}
}
{
    \eR{[English source not present in the checked-in Bencraft/Joly files.]}
}
{
}

\Verse{197}
{
    \zR{应嘉拍手道：“奇! 我在家听见说老亲翁有个衔玉生的爱子，名叫宝玉，因与小儿同名，心中甚为罕异。}
}
{
    \eR{[English source not present in the checked-in Bencraft/Joly files.]}
}
{
}

\Verse{198}
{
    \zR{后来想着这个也是常有的事，不在意了。}
}
{
    \eR{[English source not present in the checked-in Bencraft/Joly files.]}
}
{
}

\Verse{199}
{
    \zR{岂知今日一见，不但面貌相同，且举止一 般，这更奇了。”}
}
{
    \eR{[English source not present in the checked-in Bencraft/Joly files.]}
}
{
}

\Verse{200}
{
    \zR{问起年纪，“比这里的哥儿略小一岁。”}
}
{
    \eR{[English source not present in the checked-in Bencraft/Joly files.]}
}
{
}

\Verse{201}
{
    \zR{贾政便又提起承荐包勇， 问及“令郎哥儿与小儿同名”的话述了一遍。}
}
{
    \eR{[English source not present in the checked-in Bencraft/Joly files.]}
}
{
}

\Verse{202}
{
    \zR{应嘉因属意宝玉，也不暇问及那包勇 的好歹，只连连的称道：“真真罕异！”}
}
{
    \eR{[English source not present in the checked-in Bencraft/Joly files.]}
}
{
}

\Verse{203}
{
    \zR{因又拉着宝玉的手，极致殷勤。}
}
{
    \eR{[English source not present in the checked-in Bencraft/Joly files.]}
}
{
}

\Verse{204}
{
    \zR{又恐安国 公起身甚速，急须预备长行，勉强分手徐行。}
}
{
    \eR{[English source not present in the checked-in Bencraft/Joly files.]}
}
{
}

\Verse{205}
{
    \zR{贾琏宝玉送出，一路又问了宝玉好些， 然后才登车而去。}
}
{
    \eR{[English source not present in the checked-in Bencraft/Joly files.]}
}
{
}

\Verse{206}
{
    \zR{那贾琏宝玉回来见了贾政，便将应嘉问的话回了一遍。}
}
{
    \eR{[English source not present in the checked-in Bencraft/Joly files.]}
}
{
}

\Verse{207}
{
    \zR{贾政命他 二人散去。}
}
{
    \eR{[English source not present in the checked-in Bencraft/Joly files.]}
}
{
}

\Verse{208}
{
    \zR{贾琏又去张罗，算明凤姐丧事的账目。}
}
{
    \eR{[English source not present in the checked-in Bencraft/Joly files.]}
}
{
}

\Verse{209}
{
    \zR{宝玉回到自己房中，告诉了宝钗，说是：“常提的甄宝玉，我想一见不能，今 日倒先见了他父亲了。}
}
{
    \eR{[English source not present in the checked-in Bencraft/Joly files.]}
}
{
}

\Verse{210}
{
    \zR{我还听得说，宝玉也不日要到京了，要求拜望我们老爷呢。}
}
{
    \eR{[English source not present in the checked-in Bencraft/Joly files.]}
}
{
}

\Verse{211}
{
    \zR{他也说和我一模一样的，我只不信。}
}
{
    \eR{[English source not present in the checked-in Bencraft/Joly files.]}
}
{
}

\Verse{212}
{
    \zR{若是他后儿到了咱们这里来，你们都去瞧瞧， 看他果然和我像不像？”}
}
{
    \eR{[English source not present in the checked-in Bencraft/Joly files.]}
}
{
}

\Verse{213}
{
    \zR{宝钗听了道：“嗳，你说话怎么越发没前后了？}
}
{
    \eR{[English source not present in the checked-in Bencraft/Joly files.]}
}
{
}

\Verse{214}
{
    \zR{什么男人 同你一样都说出来了，还叫我们瞧去呢。”}
}
{
    \eR{[English source not present in the checked-in Bencraft/Joly files.]}
}
{
}

\Verse{215}
{
    \zR{宝玉听了，知是失言，脸上一红，连忙 的还要解说。}
}
{
    \eR{[English source not present in the checked-in Bencraft/Joly files.]}
}
{
}

\Verse{216}
{
    \zR{不知何话，下回分解。}
}
{
    \eR{[English source not present in the checked-in Bencraft/Joly files.]}
}
{
}

\Verse{217}
{
    \zR{(本章完)}
}
{
    \eR{[English source not present in the checked-in Bencraft/Joly files.]}
}
{
}
